\section{Introduction}
\label{sec:intro}

The unprecedented scale of contemporary deep neural networks has spurred a paradigm shift toward modular and parameter-efficient learning. Rather than training monolithic, multi-task systems that suffer from catastrophic forgetting and optimization interference, the modern consensus favors training specialized expert modules (e.g., parameter-efficient fine-tuning (PEFT) weights, task-specific adapters, or distinct low-rank modules) that are subsequent merged or aggregated \cite{houlsby2019parameter, hu2021lora, he2021towards}. 

In this landscape, model merging has emerged as a key technology for consolidating the knowledge of disparate specialized experts into a single unified network structure. Initial efforts focused on \emph{static model merging} through parameter-space averaging, including Weight Averaging (WA) \cite{mcmahan2017communication, izmailov2018averaging}, Model Soups \cite{wortsman2022model}, Task Arithmetic \cite{ilharco2022editing}, Fisher-Weighted Averaging \cite{matena2022merging}, RegMean \cite{jin2022dataless}, TIES-Merging \cite{yadav2023ties}, and DARE \cite{yu2023language}. While static merging effectively eliminates separate storage overheads, it invariably imposes a severe compromise: the merged model parameters are static, representing a fixed compromise across all tasks, which frequently leads to destructive parameter interference and performance degradation under high task-heterogeneity.

To resolve these limitations, \emph{dynamic model merging} has been proposed. Dynamic merging seeks to route and blend expert parameters \emph{on-the-fly} on a per-sample or per-layer basis, adjusting the parameter blending weights $\alpha(x)$ conditioned on the latent representations of each incoming input $x$. This localized parameter assembly draws inspiration from classical Mixture-of-Experts (MoE) routing \cite{jacobs1991adaptive, jordan1994hierarchical} and modern sparsely-gated layers \cite{shazeer2017outrageously, fedus2022switch}. 

Despite its conceptual appeal, existing dynamic merging methodologies are constrained by two critical and unaddressed challenges:

\begin{figure*}[t]
  \centering
  \begin{subfigure}[b]{0.32\textwidth}
    \centering
    \includegraphics[width=\textwidth]{../results/fig1_scoreboard_comparison.png}
    \caption{Scoreboard Comparison}
    \label{fig:scoreboard}
  \end{subfigure}
  \hfill
  \begin{subfigure}[b]{0.32\textwidth}
    \centering
    \includegraphics[width=\textwidth]{../results/fig2_stream_heterogeneity_audit.png}
    \caption{Heterogeneity Stream Audit}
    \label{fig:audit}
  \end{subfigure}
  \hfill
  \begin{subfigure}[b]{0.32\textwidth}
    \centering
    \includegraphics[width=\textwidth]{../results/fig3_uncertainty_mapping.png}
    \caption{Uncertainty Variance Profile}
    \label{fig:uncertainty}
  \end{subfigure}
  \caption{Empirical verification of Gaussian Process Dynamic Routing (GP-DR) on the Isolating Coordinate Sandbox. (a) Highlights the Overfitting-Optimizer Paradox, where parametric linear models collapse at test-time under data scarcity while GP-DR remains extremely stable. (b) Visualizes the dramatic recovery of GP-DR and PFSR under mixed-task streams using Micro-Batch Homogenization (MBH). (c) Illustrates exact closed-form GP posterior variance mapping, correctly identifying high epistemic uncertainty for noisy in-distribution inputs (SVHN) and true unseen out-of-distribution noise.}
  \label{fig:empirical_triad}
\end{figure*}

\begin{enumerate}
    \item \textbf{The Overfitting-Optimizer Paradox:} Dynamic routers typically employ parametric layers (such as linear projections or small multi-layer perceptrons) calibrated via backpropagation on a small set of validation data (e.g., $N=64$ samples) \cite{rostova2025tsar}. Under such extreme calibration data scarcity, gradient-descent optimizers easily memorize spurious representation-space noise. Consequently, while parametric routers achieve near-perfect training performance, their out-of-distribution and test-time performance collapses catastrophically (e.g., from an 82\% training joint score down to a 30\% test joint mean). Standard regularization methods (such as weight decay or basic dropout) are insufficient to prevent this localized memorization.
    \item \textbf{Heterogeneity Stream Collapse:} In real-world inference deployments, test inputs arrive in streaming batches containing highly heterogeneous samples from multiple distinct tasks. Because standard deep networks process inputs collectively in vectorized batches, standard routers that average features or apply uniform scaling across the batch size suffer from \emph{vectorization collapse} \cite{vance2025vrrouter}. In the presence of diverse mixed-task representations, dynamic routing coefficients average out to a flat, uniform vector across the batch size, forcing the dynamic router to collapse back to static uniform-merging levels of performance.
\end{enumerate}

To overcome these dual limitations, we reject the paradigm of parametric gradient-based routing calibration and introduce a non-parametric Bayesian alternative. In this work, we propose \textbf{Gaussian Process Dynamic Routing (GP-DR)}, a mathematically rigorous framework that completely eliminates trainable routing parameters. By treating the small pool of frozen calibration samples as spatial landmarks on the latent representation manifold, we place a Gaussian Process prior on the dynamic parameter routing mapping. Under a smooth, positive-definite Mercer kernel (such as the Radial Basis Function kernel), we solve the dynamic blending coefficients analytically as a \emph{closed-form posterior mean} in a single, stable forward matrix pass. 

Furthermore, GP-DR provides a theoretically grounded starting point for out-of-distribution (OOD) detection. While the closed-form posterior variance $\sigma^2(\psi_*)$ theoretically acts as a mathematically bounded measure of epistemic uncertainty to trigger a uniform fallback, we explicitly expose a critical \emph{unit-sphere variance collapse limitation}. Because the calibration landmarks are distributed on the compact unit sphere, random OOD noise projected onto the sphere is geometrically close to at least one landmark, rendering the GP variance blind to realistic unit-sphere OOD noise. Additionally, we empirically show that simple, unregularized distance-based heuristics (such as nearest-neighbor distances) outperform the GPR posterior variance by a wide margin under representational overlap. We openly document and analyze these limitations, providing the community with a highly transparent, mathematically rigorous account of GPR behavior on low-dimensional projection spaces.

To resolve stream-level heterogeneity, we pair GP-DR with \textbf{Micro-Batch Homogenization (MBH)}. MBH operates at the streaming buffer level, dynamically partitioning heterogeneous incoming batches into homogeneous micro-batches before they enter the modular network. This completely prevents representation-averaging collapse and restores the localized routing capabilities of GP-DR under mixed-task streams.

We conduct a rigorous empirical evaluation on a $14$-layer, $192$-dimensional Isolating Coordinate Sandbox containing four distinct tasks (MNIST, FashionMNIST, CIFAR-10, SVHN). Our results, summarized in Figure~\ref{fig:empirical_triad}, verify that GP-DR completely bypasses the Overfitting-Optimizer Paradox, achieving $72.40\%$ Joint Mean test accuracy with \emph{zero training loops}---an absolute improvement of $+42.40\%$ over standard unregularized parametric linear baselines. When deployed in heterogeneous streaming scenarios, the integration of MBH recovers GP-DR performance to $70.20\%$ ($+42.80\%$ recovery margin).

\textbf{Summary of Contributions:}
\begin{itemize}
    \item We mathematically characterize and empirically expose the \emph{Overfitting-Optimizer Paradox} and \emph{Heterogeneity Stream Collapse} in low-data dynamic model merging.
    \item We propose \textbf{Gaussian Process Dynamic Routing (GP-DR)}, a non-parametric Bayesian dynamic routing framework requiring zero trainable parameters and zero gradient updates.
    \item We derive analytical, closed-form predictive posterior mean formulas for dynamic routing and closed-form predictive posterior variance equations. We explicitly analyze their theoretical limitations, disclosing the unit-sphere variance collapse and showing where simpler distance-based heuristics outperform GPR variance under representational coupling.
    \item We demonstrate that integrating stream-level \textbf{Micro-Batch Homogenization (MBH)} completely resolves vectorization collapse under mixed-task streams.
    \item Exhaustive comparative sweeps show that GP-DR achieves exceptional stability, matching or outperforming complex, over-parameterized baselines while providing solid, provable mathematical guarantees.
\end{itemize}